\documentclass[../DoAn.tex]{subfiles}
\begin{document}
\section{Đặt vấn đề}
\label{section:1.1}

Hoạt động tư vấn pháp luật đóng vai trò cốt lõi trong việc cung cấp hỗ trợ pháp lý, giải đáp các vướng mắc và hướng dẫn người dân tuân thủ đúng quy định, nhằm bảo vệ tối đa quyền và lợi ích hợp pháp của các cá nhân cũng như tổ chức.

Trong hệ thống pháp lý, Luật Hôn nhân và Gia đình là một lĩnh vực chuyên ngành đặc thù, điều chỉnh trực tiếp các mối quan hệ phát sinh từ hôn nhân và gia đình. Phân hệ pháp luật này bao gồm toàn bộ các quy định chặt chẽ về việc kết hôn, ly hôn, quyền và nghĩa vụ của vợ chồng, quyền nuôi con, phân chia khối tài sản chung và các vấn đề khác liên quan đến gia đình. Mục tiêu của Luật Hôn nhân và Gia đình là bảo vệ quyền lợi hợp pháp của các bên trong mối quan hệ hôn nhân và gia đình, đồng thời đảm bảo sự công bằng và ổn định trong xã hội.

Hiện nay, nhu cầu nhận tư vấn chuyên sâu về lĩnh vực hôn nhân và gia đình đang tăng trưởng mạnh mẽ do các tác động xã hội phức tạp, cụ thể bao gồm các nguyên nhân như tỷ lệ ly hôn tăng và trẻ hóa, tranh chấp tài sản và quyền nuôi con, sự thay đổi trong nhận thức về bình đẳng giới; xu hướng chủ động lập thỏa thuận tài sản trước hôn nhân; cùng các vấn đề phát sinh mới như kết hôn có yếu tố nước ngoài v.v... Do mỗi quyết định dựa trên tư vấn pháp lý đều ảnh hưởng trực tiếp đến quyền lợi thân nhân hoặc tài sản lớn của người dùng nên khi nhận tư vấn, người dùng thường yêu cầu các câu trả lời phải được lập luận dựa trên các căn cứ pháp lý minh bạch được trích dẫn chính xác cụ thể với số hiệu điều khoản rõ ràng để người dùng có thể kiểm chứng, đối chiếu với các văn bản pháp luật gốc và tin tưởng vào tính chính xác của câu trả lời.

Không dừng lại ở đó, do bản chất hệ thống văn bản pháp luật về hôn nhân và gia đình tại Việt Nam vốn vô cùng phức tạp với nhiều tầng lớp đan xen giữa Luật, Nghị định, Thông tư và liên tục được cập nhật, sửa đổi nên đối với một người dùng không chuyên, việc tự xác định xem văn bản nào còn hiệu lực, văn bản nào đã bị thay thế, văn bản nào không nên áp dụng do đã bị thay thế, văn bản nào không thuộc đúng phạm vi áp dụng hoặc có sự chồng chéo về cách diễn giải là một khó khăn lớn, dễ dẫn đến những thiệt thòi, hệ lụy đáng tiếc khi vô tình áp dụng sai luật. Chính vì vậy, người dùng thường có nhu cầu nhận cảnh báo theo từng văn bản được viện dẫn, bao gồm trạng thái còn hiệu lực, đã hết hiệu lực, sắp hết hiệu lực hoặc sắp có hiệu lực thay thế; đồng thời cần nhận diện rõ từng cặp hoặc từng nhóm căn cứ pháp lý có khả năng mâu thuẫn, chồng chéo để tránh áp dụng sai luật và ý thức được các hậu quả pháp lý khác nhau khi áp dụng các căn cứ pháp lý khác nhau cho một tình huống cụ thể.

Song song với nhóm người dùng cuối trực tiếp đặt câu hỏi, một hệ thống tư vấn pháp luật tự động còn cần phục vụ nhóm người cung cấp dịch vụ (Provider), đây không nhất thiết là người dùng trực tiếp hỏi đáp, mà là cá nhân hoặc tổ chức chịu trách nhiệm triển khai, giám sát, kiểm chứng chất lượng của hệ thống chatbot trước khi cung cấp dịch vụ tư vấn cho công chúng. Do hiện tượng ảo giác (hallucination) của trí tuệ nhân tạo, Provider có nhu cầu được quan sát và kiểm chứng rõ quá trình xử lý của hệ thống. Cụ thể, Provider cần nhìn thấy rõ tiến trình suy luận của chatbot, các bước trích xuất ý định/thực thể, ngữ cảnh pháp lý được truy xuất và lý do chatbot lựa chọn điều luật đó để hoàn toàn tin tưởng và có thể tự truy vết, kiểm tra và thực hiện các sửa đổi, bổ sung cho hệ thống khi cần thiết.

Trước các nhu cầu cấp thiết đó, việc nghiên cứu và xây dựng một hệ thống chatbot có khả năng tư vấn tự động về Luật Hôn nhân và Gia đình và có khả năng minh bạch hóa luồng suy luận trước khi đưa ra câu trả lời tư vấn mang lại những ý nghĩa thực tiễn vô cùng to lớn. Hệ thống này đóng vai trò như một công cụ bổ trợ đắc lực cho các cơ quan tư vấn công lập hiện tại, giúp giảm tải hiệu quả áp lực công việc cho đội ngũ cán bộ ngành luật; giúp người dân đại chúng nhanh chóng tiếp cận được những tư vấn pháp luật chính xác, minh bạch; giúp những người cung cấp dịch vụ chatbot (Provider) có thêm công cụ theo dõi, kiểm thử và bảo trì hệ thống chatbot, nhằm đảm bảo chất lượng câu trả lời của chatbot khi đưa ra tư vấn cho người dùng.

Mặc dù mang ý nghĩa quan trọng, việc hiện thực hóa một chatbot tư vấn pháp lý tự động như vậy vẫn phải đối mặt với những rào cản mang tính đặc thù. Hệ thống pháp luật về hôn nhân và gia đình tại Việt Nam mang tính chất phức tạp, và thường xuyên được cập nhật, gây khó khăn trong tiếp cận, tra cứu và áp dụng. Một khó khăn quan trọng khác là các văn bản pháp luật có thể tồn tại quan hệ sửa đổi, bổ sung, thay thế, hướng dẫn hoặc thậm chí có khả năng chồng chéo, mâu thuẫn trong một số tình huống áp dụng cụ thể. Vì vậy, chatbot tư vấn pháp luật hôn nhân gia đình không chỉ cần truy xuất đúng điều luật, mà còn cần kiểm tra quan hệ giữa các văn bản viện dẫn, xử lý khả năng xung đột căn cứ pháp lý, giải thích rõ vì sao các căn cứ đó được sử dụng trong câu trả lời và hiển thị luồng xử lý đủ chi tiết để Provider giám sát quá trình hệ thống vận hành.

Hiện nay, các công cụ hỗ trợ tra cứu và tư vấn pháp luật tại Việt Nam đã có những bước phát triển nhất định với sự xuất hiện của các sản phẩm như AITracuuluat, AI luật và Chatbot của Bộ Tư pháp. Bên cạnh đó, các mô hình ngôn ngữ lớn phổ biến như Gemini cũng thường xuyên được người dùng sử dụng để giải đáp các vướng mắc pháp lý. 

Qua quá trình khảo sát và đánh giá, các sản phẩm hiện có bộc lộ những ưu điểm và hạn chế khác nhau trong việc giải quyết các bài toán pháp lý. Công cụ AITracuuluat đôi khi nêu ra các căn cứ pháp lý không chính xác, không liên quan đến nhóm vấn đề pháp lý trong câu hỏi của người dùng hoặc viện dẫn căn cứ pháp lý đã hết hiệu lực. Trong khi đó, mặc dù chatbot của Bộ Tư pháp và AI Luật đảm bảo tính cập nhật và độ chính xác cao, hệ thống này lại chưa có khả năng cảnh báo các văn bản đã hết hiệu lực, sắp hết hiệu lực hoặc sắp có hiệu lực và cũng chưa giải thích lý do vì sao áp dụng các căn cứ pháp lý được nêu ra trong câu trả lời. Đối với các mô hình ngôn ngữ lớn như Gemini, câu trả lời thường xuyên thiếu sự viện dẫn đầy đủ các căn cứ pháp lý cần thiết và gặp khó khăn trong việc cập nhật thời hạn hiệu lực của văn bản quy phạm pháp luật. Các hệ thống hiện có cũng chưa cung cấp cơ chế hiển thị chi tiết luồng suy luận trước khi đưa ra câu trả lời để Provider có thể kiểm tra, giám sát quá trình xử lý, suy luận của chatbot.

Tựu trung lại, các chatbot tư vấn pháp luật Việt Nam hiện hành còn gặp phải 4 nhóm vấn đề chính: (i) Đôi khi xảy ra hiện tượng trích dẫn các căn cứ pháp lý thiếu chính xác hoặc không bao phủ đủ các nhóm vấn đề pháp lý trong câu hỏi. (ii) Chưa có cơ chế cảnh báo tình trạng hiệu lực văn bản (bao gồm văn bản đã hết hiệu lực, sắp hết hiệu lực hoặc sắp có hiệu lực), cảnh báo về các căn cứ có khả năng mâu thuẫn khi áp dụng trong các tình huống cụ thể. (iii) Chưa có cơ chế hiển thị luồng xử lý, suy luận của chatbot trước khi đưa ra câu trả lời.

Việc giải quyết bài toán tra cứu và tư vấn pháp luật tự động mang lại những giá trị thực tiễn thiết thực cho cộng đồng. Xử lý những rào cản nêu trên giúp người dân nhanh chóng tiếp cận nguồn dữ liệu pháp lý minh bạch, chính thống, có thể kiểm chứng theo từng căn cứ pháp lý; đồng thời giúp Provider có cơ sở quan sát, kiểm thử và cải tiến hệ thống theo từng bước xử lý cụ thể. Một cơ chế tư vấn luật pháp hiệu quả sẽ đóng vai trò bổ trợ đắc lực cho các hệ thống tư vấn công lập, góp phần giảm thiểu áp lực công việc cho đội ngũ cán bộ ngành luật. Khuôn mẫu giải quyết bài toán này không chỉ giới hạn trong phạm vi Luật Hôn nhân và gia đình mà còn sở hữu tiềm năng mở rộng sang nhiều phân hệ pháp luật chuyên ngành khác.
\section{Mục tiêu và phạm vi đề tài}
\label{section:1.2}
Từ những đánh giá, phân tích về các vấn đề, tồn tại trên, ĐATN hướng tới mục tiêu xây dựng một hệ thống chatbot tư vấn Luật Hôn nhân và Gia đình có khả năng hỗ trợ hai nhóm người dùng chính: Thứ nhất là người dùng cuối cần nhận câu trả lời được lập luận dựa trên căn cứ pháp lý chính xác, minh bạch, có cảnh báo hiệu lực văn bản và mâu thuẫn giữa các căn cứ pháp lý (nếu có); Thứ hai là những người cung cấp dịch vụ chatbot (Provider) với nhu cầu cần quan sát được luồng xử lý của hệ thống để kiểm thử, bảo trì và kiểm chứng chất lượng trước khi cung cấp dịch vụ tư vấn cho công chúng. Trọng tâm của đồ án là thiết lập một cơ chế truy xuất các căn cứ pháp lý chính xác, đảm bảo mọi phản hồi đều dựa trên các căn cứ pháp lý phù hợp với thời điểm áp dụng, được viện dẫn chi tiết, đồng thời có cơ chế cảnh báo hiệu lực văn bản và mâu thuẫn pháp lý, và minh bạch hóa luồng xử lý của hệ thống để Provider có thể theo dõi, kiểm tra và cải tiến hệ thống một cách hiệu quả.

Về phạm vi chủ đề, phần mềm được phát triển có khả năng trả lời các câu hỏi thuộc các nhóm nội dung liên quan đến Luật Hôn nhân và Gia đình, bao gồm: Kết hôn, Quan hệ giữa vợ và chồng trong hôn nhân, Chấm dứt hôn nhân, Chế độ tài sản của vợ chồng, Quan hệ giữa cha mẹ và con, Quan hệ hôn nhân và gia đình có yếu tố nước ngoài. (Chi tiết về các dạng câu hỏi con thuộc các nhóm nội dung lớn này sẽ được trình bày tại chương \ref{chapter:Experiment}) và chương \ref{chapter:Methodology}). Trong các chủ đề này, chatbot tập trung hỗ trợ hai dạng câu hỏi chính. Dạng thứ nhất là câu hỏi lý thuyết pháp luật, tức các câu hỏi về nội dung điều luật, khái niệm, nguyên tắc pháp luật, điều kiện pháp lý, quyền và nghĩa vụ của các bên hoặc hệ quả pháp lý được quy định trong văn bản pháp luật. Dạng thứ hai là câu hỏi tình huống pháp luật, tức các câu hỏi mô tả một sự kiện cụ thể và yêu cầu hệ thống phối hợp áp dụng một hoặc nhiều điều luật để đưa ra nhận định tư vấn. 

Do giới hạn về thời gian và khả năng, Chatbot ĐATN phát triển hiện tại chưa bảo đảm chất lượng câu trả lời đối với các câu hỏi tư vấn liên quan đến thủ tục hành chính như thành phần hồ sơ, giấy tờ, trình tự thực hiện, cơ quan tiếp nhận hoặc biểu mẫu cụ thể. Hệ thống cũng chưa giải đáp được các câu hỏi cần kết hợp căn cứ từ luật tố tụng dân sự và thủ tục tố tụng dân sự liên quan đến hôn nhân và gia đình, chẳng hạn thủ tục ly hôn, trình tự nộp đơn, thẩm quyền tòa án, các bước tham gia, xét xử vụ án tại phiên tòa hoặc đưa ra các tư vấn về phán quyết có thể của tòa đối với một vụ án.

Về tiêu chí đánh giá câu trả lời, một câu trả lời lý thuyết được xem là đạt yêu cầu khi nêu đúng nội dung pháp luật được hỏi, trích dẫn chính xác văn bản và chi tiết đến từng điều, khoản, điểm, đồng thời đi kèm trạng thái hoặc thời hạn hiệu lực của từng căn cứ pháp lý. Một câu trả lời tình huống được xem là đạt yêu cầu khi xác định đúng nhóm vấn đề pháp lý trong câu hỏi (Dạng câu hỏi pháp lý), nhận diện đúng các yếu tố pháp lý quan trọng của tình huống, truy xuất đủ các căn cứ chính cần áp dụng. Trong mọi trường hợp, các căn cứ pháp lý được viện dẫn phải là căn cứ còn hiệu lực tại thời điểm truy vấn hoặc có hiệu lực tại thời điểm xảy ra sự kiện pháp lý trong câu hỏi. Đối với các trường hợp câu hỏi mà câu trả lời phải sử dụng căn cứ đã hết hiệu lực, sắp hết hiệu lực hoặc sắp bị thay thế thì hệ thống phải cảnh báo rõ cho người dùng về trạng thái đó và giải thích lý do vì sao cần viện dẫn đến căn cứ đó. Nếu tồn tại các căn cứ pháp lý có khả năng chồng chéo, mâu thuẫn khi áp dụng, hệ thống phải trích dẫn chính xác các căn cứ liên quan và cảnh báo người dùng về khả năng phát sinh mâu thuẫn trong cách áp dụng. Một câu trả lời không được xem là đạt yêu cầu nếu áp dụng các điều luật không chính xác hoặc không liên quan đến tình huống, viện dẫn sai điều luật hoặc áp dụng các căn cứ không phù hợp về hiệu lực.

Đối với tiêu chí minh bạch hóa luồng hoạt động của hệ thống, giao diện cần hiển thị rõ đầu vào và đầu ra của các tác nhân (agent) được kích hoạt trong quá trình trả lời, bao gồm bước phân loại câu hỏi và lựa chọn công cụ truy xuất (router agent), trích xuất ý định/thực thể và truy xuất ngữ cảnh pháp lý (retriever agent). Bên cạnh đó, Provider cần xem được chính xác toàn bộ các node và relationship trong đồ thị tri thức đã được truy vấn trong quá trình truy xuất và tổng hợp ngữ cảnh pháp lý cho LLM. Phần minh bạch hóa luồng hoạt động của hệ thống được xem là đạt yêu cầu khi Provider có thể truy vết từ câu hỏi đầu vào, qua các agent và dữ liệu đồ thị được truy xuất, đến câu trả lời cuối cùng của chatbot.

\section{Định hướng giải pháp}
\label{section:1.3}

ĐATN đã khảo sát, phân tích, nêu ra vấn đề của các ứng dụng chatbot tư vấn pháp luật hiện hành tại Việt Nam cùng những hạn chế liên quan đến việc trích dẫn các căn cứ pháp lý chính xác là cơ sở để trả lời cho câu hỏi của người dùng, thiếu cơ chế cảnh báo hiệu lực văn bản và xử lý các tình huống mâu thuẫn pháp lý thực tế. Để đáp ứng, hoàn thiện, khắc phục những hạn chế, tồn tại đã chỉ ra ở trên, ĐATN đưa ra phương án giải quyết cho từng vấn đề: (i) mô hình hóa các thực thể pháp lý và mối liên hệ giữa chúng (ví dụ: giữa “quyền nuôi con” và “độ tuổi trẻ em” hay “thỏa thuận của cha mẹ”) và các mối quan hệ hiệu lực, sửa đổi, hướng dẫn bổ sung, thay thế, bác bỏ giữa các văn bản quy phạm pháp luật của Luật Hôn nhân và Gia đình Việt Nam dưới dạng đồ thị tri thức (Knowledge Graph); (ii) Thiết kế cơ chế tự động kiểm tra và cảnh báo tình trạng hiệu lực văn bản pháp lý (iii) Thiết kế cơ chế tự động kiểm tra và cảnh báo tình trạng hiệu lực, đồng thời phát hiện và xử lý mâu thuẫn trong các căn cứ pháp lý (nếu có).

Từ đầu vào là một câu hỏi của người dùng bằng tiếng Việt tự nhiên thuộc các chủ đề của Luật Hôn nhân và Gia đình , hệ thống đề xuất cần đưa ra đầu ra bao gồm một câu trả lời hoàn chỉnh đính kèm trích dẫn điều khoản chính xác, hiển thị tường minh luồng xử lý của các agent (agent trace) và ngữ cảnh pháp lý được truy xuất từ đồ thị tri thức. Trước hết, kết quả phản hồi này cần thỏa mãn ràng buộc về tính chính xác pháp lý tuyệt đối, đảm bảo các điều luật được viện dẫn đều là văn bản đang còn hiệu lực hoặc hoàn toàn phù hợp với mốc thời gian và bối cảnh của tình huống được hỏi. Ngoài ra, trong câu trả lời cần có những cảnh báo pháp lý khi áp dụng các căn cứ pháp lý đã hết hiệu lực hoặc sắp hết hiệu lực, giải thích tại sao lại áp dụng các căn cứ pháp lý đó trong trường hợp có mâu thuẫn, chồng chéo trong những căn cứ pháp lý.

Quy trình xử lý của kiến trúc đề xuất bao gồm 5 bước chính, cụ thể thực hiện theo thứ tự như sau. (i) Ban đầu, sử dụng mô hình ngôn ngữ lớn kết hợp Prompt Engineering để phân loại dạng câu hỏi, xác định (những) chủ đề pháp lý liên quan đến câu hỏi người dùng dùng. (ii) Tiếp theo, điều phối các Agent Retriever chuyên để xử lý cho từng dạng câu hỏi, các agent có nhiệm vụ phải trích xuất ra các ý định (intent) và thực thể (entity) trong câu hỏi của người dùng, sau đó điền các tham số trích xuất được vào mẫu câu lệnh cypher (cypher template) để sinh ra câu lệnh Cypher để truy xuất các căn cứ pháp lý cần để trả lời cho câu hỏi của người dùng từ đồ thị tri thức. (iii) Sau đó tiến hành lọc các căn cứ còn hiệu lực (hoặc phù hợp với thời điểm xảy ra sự kiện pháp lý trong câu hỏi) và các căn cứ có khả năng mâu thuẫn (iv) Tiếp theo, sắp xếp các căn cứ pháp lý thành các nhóm là căn cứ chính, căn cứ tham chiếu bổ trợ, căn cứ hướng dẫn, căn cứ thay thế và căn cứ mâu thuẫn. (v) Cuối cùng, gọi API mô hình ngôn ngữ lớn để sinh câu trả lời cuối cùng đảm bảo có trích dẫn căn cứ pháp lý chi tiết, cảnh báo hiệu lực văn bản và giải thích ngắn gọn mối quan hệ giữa các căn cứ pháp lý, lí do chọn các căn cứ pháp lý đó. 

Công nghệ được sử dụng cho ĐATN bao gồm nhiều thành phần khác nhau. Thư viện Chainlit hỗ trợ việc lập trình giao diện Web Chatbot UI một cách nhanh chóng, cung cấp sẵn các cơ chế quản lý session hỏi đáp và hiển thị trực quan luồng agent trace cho người dùng. Ngôn ngữ lập trình Python được chọn làm ngôn ngữ chủ đạo nhờ sở hữu hệ sinh thái thư viện phong phú, tối ưu cho các bài toán xử lý ngôn ngữ tự nhiên và khoa học dữ liệu. Cơ sở dữ liệu đồ thị Neo4j cho phép mô hình hóa, lưu trữ và truy vấn nhanh chóng các mối quan hệ thực thể pháp lý, kết hợp với PostgreSQL nhằm quản lý thông tin người dùng và lưu trữ lịch sử hội thoại. Để thực hiện các tác vụ suy luận và trích xuất dữ liệu có cấu trúc, hệ thống tích hợp API của các mô hình ngôn ngữ lớn tiên tiến bao gồm dòng GPT-4.1, gpt 3o và gpt-4o, Gemini 3.1  nhằm đảm bảo chất lượng sinh ngôn ngữ. Cuối cùng, thư viện RAGAS được ứng dụng làm nền tảng đánh giá định lượng, giúp đo lường chính xác hiệu năng của toàn bộ luồng truy xuất.

Các đóng góp chính của đồ án bao gồm: (i) Xây dựng đồ thị tri thức cho Luật hôn nhân gia đình 2014 và và mẫu truy vấn Cypher (Cypher template) cho từng dạng câu hỏi (ii) Cơ chế kiểm tra tình trạng hiệu lực, xử lý phân loại các văn bản
pháp luật (căn cứ chính, căn cứ bổ trợ, căn cứ hướng dẫn, sửa đổi, thay thế) và các căn cứ chồng chéo, xung đột (iii) thiết kế và xây dựng bộ Benchmark dataset tiêu chuẩn gồm tối thiểu 600 câu hỏi - đáp pháp lý có gán nhãn cấu trúc chi tiết phục vụ kiểm thử.

\section{Bố cục đồ án}
\label{section:1.4}

Phần còn lại của báo cáo đồ án tốt nghiệp này được tổ chức như sau.

Chương \ref{chapter:Related_works} trình bày tổng quan kết quả khảo sát các hệ thống chatbot tư vấn pháp luật hiện hành và thực trạng áp dụng trợ lý ảo để tư vấn luật pháp nói chung và luật hôn nhân gia đình nói riêng tại Việt Nam. Căn cứ vào các ưu nhược điểm của những công cụ đang có trên thị trường, các giải pháp hiện hành và lấy đó làm cơ sở để xác định rõ bài toán mà đồ án cần tập trung giải quyết.

Trong Chương \ref{chapter:Methodology}, đồ án giới thiệu các cơ sở lý thuyết và nền tảng công nghệ lõi được ứng dụng để phát triển hệ thống phần mềm. Nội dung tập trung làm rõ các khái niệm về Đồ thị tri thức (Knowledge Graph), kỹ thuật Retrieval-Augmented Generation (RAG) và kiến trúc Agentic AI. Bên cạnh đó, chương này cũng đề cập đến các công cụ lưu trữ dữ liệu đồ thị như Neo4j và nền tảng giao diện Chainlit.

Chương \ref{chapter:Experiment} cung cấp nội dung chi tiết về quá trình thiết kế hệ thống và quá trình thử nghiệm hệ thống và phân tích các kết quả đánh giá hiệu năng. Dựa trên bộ dữ liệu kiểm chuẩn chuẩn hóa (benchmark dataset), chương này trình bày các số liệu thực nghiệm liên quan đến độ chính xác trong việc lựa chọn công cụ của tác nhân AI, khả năng truy xuất ngữ cảnh và chất lượng câu trả lời cuối cùng. Các kết quả này được đối chiếu trực tiếp với hệ thống chỉ tiêu đề ra thông qua bộ chỉ số RAGAS và các tiêu chí xử lý văn bản luật khác.

Chương \ref{chapter:SolutionAndContribution} trình bày đóng góp cốt lõi của đồ án thông qua việc mô tả chi tiết giải pháp kỹ thuật và quá trình xây dựng hệ thống Agentic GraphRAG. Chương này mô tả việc thiết kế đồ thị knowledge graph để mô hình hóa các thực thể pháp lý cùng mối quan hệ giữa Luật hôn nhân và gia đình năm 2014 và các văn bản dưới luật liên qua và cơ chế tự động kiểm tra tình trạng hiệu lực văn bản, áp dụng các văn bản theo đúng cấp bậc các văn bản pháp luật và giải quyết mâu thuẫn giữa các điều luật (nếu có). Cuối cùng, trình bày cách xây dựng tập benchmark dataset gồm 600 câu hỏi, bao phủ các dạng câu hỏi chính liên quan đến Luật hôn nhân và gia đình.

Cuối cùng, Chương \ref{chapter:conclusion} tổng kết lại các kết quả đã đạt được của quá trình nghiên cứu và phát triển phần mềm. Bên cạnh việc khẳng định mức độ hoàn thành các mục tiêu ban đầu, chương này cũng đưa ra những nhận định khách quan về các mặt còn hạn chế của hệ thống hiện tại, từ đó đề xuất các hướng phát triển và mở rộng tiềm năng của công cụ tư vấn pháp lý trong tương lai.

\end{document}
